\section{Centered norms: twenty-four and twenty-two integers}
\label{sec:centered}

Two additional changes improve the variable/degree tradeoff. Both use the
integral coordinate shear
\begin{equation}
 y_{\rm old}=fo+y,\qquad y=y_{\rm old}-fo.
 \label{eq:auxshear}
\end{equation}
This is a bijection on integer coordinates, since $f,o$ are retained.
Together with $V=fo-c$, it gives the exact factorization
\begin{equation}
 y_{\rm old}^{\,2}-V^2=(c+y)(2fo+y-c).
 \label{eq:auxdiff}
\end{equation}
The leading square cancels, saving one degree in this difference.
These equations use $y$ for the centered root; it no longer denotes
the uncentered $y$ of Section~\ref{sec:integer}.

\subsection{A degree-ten equation in twenty-four integers}

Start with the 26-integer octic described above. Retain $\zeta$ as an
unknown, set $k=\eta+\zeta$, eliminate $D=X+ac+\gamma H$ and $V=fo-c$,
and apply \eqref{eq:auxshear}. The centered auxiliary residual is
\begin{equation}
 A_{\rm aux}=(fo+y)^2-I^2(c+y)(2fo+y-c)-1.
 \label{eq:aux24}
\end{equation}
Its degree is five. The directly substituted main norm
$N=(X+ac+\gamma H)^2-\Delta c^2-1$ still has degree six. A different
row representing that norm will lower its degree.

Put $b=ac+\gamma H$ and $Y=a-E$, as before, and define
\begin{equation}
 M=XE+2Eb+Y\bigl[H(2ac\gamma+H\gamma^2-c^2)-1\bigr].
 \label{eq:main24}
\end{equation}
Since $\Delta=a^2+H$, direct expansion gives
\[
 b^2-\Delta c^2=H(2ac\gamma+H\gamma^2-c^2).
\]
Consequently the following identity holds as an identity of polynomials:
\begin{equation}
 \boxed{M=YN+(X+2b)(E-XY).}
 \label{eq:nonzeroscale}
\end{equation}
The retained row $E-XY=0$ therefore makes $M=YN$. The other retained
row gives $Y=(2v+1)q^3$. This is never zero at an integer solution:
$2v+1\ne0$, and $q=\beta J+1\ne0$ because $\beta>1$ and $J$ is an
integer. Thus replacing $N=0$ by $M=0$ is reversible on the common
zeros of the other rows. This explicit nonvanishing argument excludes
the extra zero component that an arbitrary multiplication could introduce.
No positivity, native calibration or Pell theorem is used to establish it.

The unknowns of $P_{24}$ are the twenty-three names in
\eqref{eq:vars23}, together with $\zeta$; $y$ has the centered meaning.
Keep the definitions of $q,C,Y,H,\Delta,X,R,G,\Pi$ from
Section~\ref{sec:integer}, but now use $k=\eta+\zeta$, $V=fo-c$,
and $b=ac+\gamma H$. Its twelve residuals are
\begin{align}
 s_1&=Y-(2v+1)q^3, &
 s_2&=E-XY,\nonumber\\
 s_3&=c-\eta-kY, &
 s_4&=\tau^2+(2\tau-k)E(c-\eta)-1,\nonumber\\
 s_5&=M, &
 s_6&=k-R-1-hE,\nonumber\\
 s_7&=(K_0q+X)C-qF-Z(q-1), &
 s_8&=k+V-cj-hE-1,\nonumber\\
 s_9&=I-(2r+1)c^2, &
 s_{10}&=I^2-\Delta(f^2-1),\nonumber\\
 s_{11}&=A_{\rm aux}, &
 s_{12}&=G\eta\zeta-(t^2-G-\eta\zeta+1)^2-\Pi.
 \label{eq:24residuals}
\end{align}
The single equation is
\begin{equation}
 \boxed{P_{24}=\sum_{\nu=1}^{12}s_\nu^2=0,
 \qquad\mathbf x\in\Z^{24}.}
 \label{eq:P24}
\end{equation}
Every residual has degree at most five. The complete degree-ten part is
\begin{equation}
 \boxed{\lead_{10}(P_{24})=
 64a^4\gamma^2(c+2\gamma)^2(a-E)^2
 +4I^4f^2o^2(c+y)^2.}
 \label{eq:leading24}
\end{equation}
These are the leading squares of $M$ and $A_{\rm aux}$. In particular,
$[I^4f^2o^2y^2]P_{24}=4$, independently of the fixed ports.
The supplied circuit has 135 binary arithmetic operations.

\begin{theorem}[The 24-integer reduction]
The polynomial $P_{24}$ has exactly twenty-four unrestricted integer
unknowns and exact degree ten. Its integer solvability is equivalent to
that of the 26-integer octic and hence to the pinned finite-refutation
predicate. Its unsolvability satisfies \eqref{eq:target} for
$\mathcal D=\Z^{24}$.
\end{theorem}
\begin{proof}
The monic definitions of $D,V$ have polynomial inverse restorations.
The shear \eqref{eq:auxshear} is invertible over the integers.
Finally \eqref{eq:nonzeroscale}, together with $E=XY$ and the proved
$Y\ne0$, makes the main-norm replacement reversible. Each step preserves
existence of common integer zeros. A sum of integer squares vanishes
exactly when all its residuals vanish. The degree certificate
\eqref{eq:leading24} proves the remaining claim.
\end{proof}

\subsection{A degree-fourteen equation in twenty-two integers}

Start with $P_{23}$ rather than $P_{24}$, so that $\zeta$ is again the
abbreviation $1-V+jc+hE-\eta$. Apply the same shear to $y$. Define
\[
 Q=\Delta(f^2-1),\qquad S=I^2-Q,
 \qquad A_{\rm old}=y_{\rm old}^{\,2}
       -I^2(y_{\rm old}^{\,2}-V^2)-1.
\]
The residual $S$ is already present. Replace $A_{\rm old}$ by
\begin{equation}
 A_{\rm new}=(fo+y)^2-Q(c+y)(2fo+y-c)-1.
 \label{eq:aux22}
\end{equation}
The exact identity
\begin{equation}
 A_{\rm new}=A_{\rm old}+S(y_{\rm old}^{\,2}-V^2)
 \label{eq:row22}
\end{equation}
shows that this is an invertible elementary row operation in the presence
of $S=0$. It needs no division or nonvanishing assumption.

Now $I$ occurs only in its definition $I=(2r+1)c^2$ and in $S$.
Eliminate it, deleting its definition row. The strong row becomes
\begin{equation}
 S_{22}=(2r+1)^2c^4-\Delta(f^2-1).
 \label{eq:strong22}
\end{equation}
The complete single equation can therefore be written using the already
displayed $r_1,\ldots,r_7,r_{11}$ as
\begin{equation}
 \boxed{P_{22}=\sum_{\nu=1}^{7}r_\nu^2
              +S_{22}^{,2}+A_{\rm new}^{,2}+r_{11}^{,2}=0,
 \qquad\mathbf x\in\Z^{22}.}
 \label{eq:P22}
\end{equation}
The twenty-two unknowns are precisely \eqref{eq:vars23} with $I$
removed. In the unchanged residuals and abbreviations, $V=fo-c$ and
none depends on $y$, so the centered meaning of $y$ creates no further
replacement. The fixed coefficient tuple is unchanged.

The only degree-seven residual is $A_{\rm new}$. Its leading form is
$-2a^2f^3o(c+y)$; all other residuals have degree at most six. Hence
\begin{equation}
 \boxed{\lead_{14}(P_{22})=4a^4f^6o^2(c+y)^2,
 \qquad[a^4f^6o^2y^2]P_{22}=4.}
 \label{eq:leading22}
\end{equation}
The degree is exactly fourteen at the actual fixed coefficients.
The supplied circuit uses 123 binary operations.

\begin{theorem}[The 22-integer reduction]
The polynomial $P_{22}$ has twenty-two unrestricted integer unknowns,
ten squared residuals, and exact degree fourteen. Its integer solvability
is equivalent to that of $P_{23}$ and thus to finite ZFC refutability
through the pinned calibration. In particular \eqref{eq:target} holds
for $(P,\mathcal D)=(P_{22},\Z^{22})$.
\end{theorem}
\begin{proof}
The root shear is an integral bijection. In the common-zero system,
\eqref{eq:row22} has the inverse row operation of subtracting
$S(y_{\rm old}^{,2}-V^2)$. After this change the monic definition of
$I$ can be eliminated, with its displayed integer polynomial as inverse
restoration. These steps preserve the entire integer zero set up to the
specified coordinate maps. The previous soundness and completeness
implications for $P_{23}$ therefore apply. Equation~\eqref{eq:leading22}
proves exact degree fourteen.
\end{proof}

Direct substitution of $I=(2r+1)c^2$ into the old auxiliary norm would
raise the degree much more. The order matters: remove $I$ from that norm
using the strong row, cancel the leading square by centering, and only
then eliminate $I$. The 22-integer result and the degree-ten companion
are distinct useful tradeoffs; neither dominates the other in both
variable count and degree.
